% 使用 XeLaTeX 编译。
\documentclass[UTF8,a4paper,11pt,fontset=windows]{ctexart}
\usepackage[margin=23mm]{geometry}
\usepackage{amsmath,amssymb,mathtools,enumitem,fancyhdr,hyperref,xcolor}
\hypersetup{hidelinks,pdftitle={拓扑学第一次作业详细答案}}
\pagestyle{fancy}
\fancyhf{}
\fancyhead[L]{拓扑学}
\fancyhead[R]{第一次作业 · 详细答案}
\fancyfoot[C]{\thepage}
\setlength{\headheight}{15pt}
\setlength{\parindent}{2em}
\setlength{\parskip}{3pt}
\setlist[enumerate]{label=(\alph*),leftmargin=2.7em,itemsep=5pt,topsep=5pt}
\newcommand{\problem}[1]{\section*{第 #1 题}}
\newcommand{\idea}{\par\noindent\textcolor{gray}{\textbf{思路。}}}
\newcommand{\solution}{\par\noindent\textbf{解答。}}
\newcommand{\R}{\mathbb R}
\begin{document}

\begin{center}
  {\LARGE\bfseries 拓扑学第一次作业：详细答案}\par\vspace{6pt}
  {\large Armstrong《Basic Topology》第 2 章}
\end{center}

\noindent\textbf{范围：}第 1、2、3、4、5、6、9、14、16、17 题。下文中，$\overline A$、$A^\circ$ 和 $\partial A$ 分别表示 $A$ 的闭包（closure）、内部（interior）和边界（frontier）。

\problem{1}
\idea 闭包适合用“包含原集合的最小闭集”来处理，内部适合用“包含于原集合的最大开集”来处理。需要说明等号不成立时，给出一个反例即可。

\solution 设 $A,B\subseteq X$。
\begin{enumerate}
  \item 因为 $A\subseteq A\cup B$ 且 $B\subseteq A\cup B$，闭包的单调性给出
  \[
    \overline A\cup\overline B\subseteq\overline{A\cup B}.
  \]
  另一方面，$\overline A$ 和 $\overline B$ 都是闭集，所以有限并 $\overline A\cup\overline B$ 仍是闭集；它又包含 $A\cup B$。由闭包的最小性，
  \[
    \overline{A\cup B}\subseteq\overline A\cup\overline B.
  \]
  两个包含关系合起来便得到
  \[
    \overline{A\cup B}=\overline A\cup\overline B.
  \]

  \item 由 $A\cap B\subseteq A$ 和 $A\cap B\subseteq B$，闭包的单调性给出
  \[
    \overline{A\cap B}\subseteq\overline A,
    \qquad
    \overline{A\cap B}\subseteq\overline B.
  \]
  因而
  \[
    \overline{A\cap B}\subseteq\overline A\cap\overline B.
  \]
  该包含关系可能是严格的。例如在通常实数直线中取
  \[
    A=(-\infty,0),\qquad B=(0,+\infty).
  \]
  此时 $A\cap B=\varnothing$，故 $\overline{A\cap B}=\varnothing$；而
  \[
    \overline A\cap\overline B=(-\infty,0]\cap[0,+\infty)=\{0\}.
  \]

  \item $\overline A$ 本身是闭集。闭集的闭包等于它本身，所以
  \[
    \overline{\overline A}=\overline A.
  \]

  \item $A^\circ\subseteq A\subseteq A\cup B$，而且 $A^\circ$ 是开集，因此 $A^\circ\subseteq(A\cup B)^\circ$。同理，$B^\circ\subseteq(A\cup B)^\circ$。于是
  \[
    (A\cup B)^\circ\supseteq A^\circ\cup B^\circ.
  \]
  等号不一定成立。仍取
  \[
    A=(-\infty,0],\qquad B=[0,+\infty).
  \]
  则 $A\cup B=\R$，所以 $(A\cup B)^\circ=\R$；但
  \[
    A^\circ\cup B^\circ=(-\infty,0)\cup(0,+\infty)=\R\setminus\{0\}.
  \]

  \item 因为 $A^\circ\cap B^\circ$ 是开集，并且包含于 $A\cap B$，所以
  \[
    A^\circ\cap B^\circ\subseteq(A\cap B)^\circ.
  \]
  反过来，$(A\cap B)^\circ$ 是同时包含于 $A$ 和 $B$ 的开集。由内部的最大性，
  \[
    (A\cap B)^\circ\subseteq A^\circ,
    \qquad
    (A\cap B)^\circ\subseteq B^\circ.
  \]
  因而 $(A\cap B)^\circ\subseteq A^\circ\cap B^\circ$。故
  \[
    (A\cap B)^\circ=A^\circ\cap B^\circ.
  \]

  \item $A^\circ$ 是开集，而开集的内部就是它本身，因此
  \[
    (A^\circ)^\circ=A^\circ.
  \]
\end{enumerate}

\problem{2}
\idea 无限多个闭集的并不一定闭。构造一列趋向某个缺失极限点的单点集即可。

\solution 对每个正整数 $n$，令
\[
  F_n=\left\{\frac1n\right\}.
\]
每个 $F_n$ 都是通常实数直线中的闭集。然而
\[
  \bigcup_{n=1}^{\infty}F_n=\left\{1,\frac12,\frac13,\ldots\right\}
\]
不是闭集，因为 $1/n\to0$，所以 $0$ 是这个并集的极限点，但 $0$ 不属于该并集。

\problem{3}
\idea 使用
\[
  \partial A=\overline A\setminus A^\circ.
\]
第三小题的关键是先求振荡曲线 $G=\{(x,\sin(1/x)):x>0\}$ 的闭包。

\solution
\begin{enumerate}
  \item 令
  \[
    A=\{(x,y)\in\R^2:1<x^2+y^2\leq2\}.
  \]
  内圆周 $x^2+y^2=1$ 不在 $A$ 中，外圆周 $x^2+y^2=2$ 虽在 $A$ 中，但两者都不是内点。于是
  \[
    A^\circ=\{(x,y):1<x^2+y^2<2\},
  \]
  \[
    \overline A=\{(x,y):1\leq x^2+y^2\leq2\},
  \]
  \[
    \partial A=\{(x,y):x^2+y^2=1\}\cup
    \{(x,y):x^2+y^2=2\}.
  \]

  \item 令 $B$ 为去掉两条坐标轴后的平面，即
  \[
    B=\{(x,y)\in\R^2:xy\neq0\}.
  \]
  它是四个开象限的并，因此是开集，故 $B^\circ=B$。坐标轴上的任一点的每个开圆盘都与某个象限相交，所以每个坐标轴上的点也属于 $\overline B$。因此
  \[
    B^\circ=B,\qquad \overline B=\R^2,
  \]
  \[
    \partial B=\{(x,y):x=0\}\cup\{(x,y):y=0\}.
  \]

  \item 记
  \[
    G=\{(x,\sin(1/x)):x>0\},\qquad C=\R^2\setminus G.
  \]
  当 $x\to0^+$ 时，$\sin(1/x)$ 在 $[-1,1]$ 中振荡。对每个 $t\in[-1,1]$，可选取趋于无穷的数列 $u_k$，使 $\sin u_k=t$；令 $x_k=1/u_k$，便有
  \[
    (x_k,\sin(1/x_k))\longrightarrow(0,t).
  \]
  所以 $\{0\}\times[-1,1]\subseteq\overline G$。反之，若 $G$ 中一点列收敛，则其第一坐标的极限不可能小于 $0$；第一坐标极限若大于 $0$，由函数 $\sin(1/x)$ 的连续性，极限仍在 $G$ 中；第一坐标极限若等于 $0$，第二坐标只能落在 $[-1,1]$ 中。因此
  \[
    \overline G=G\cup\bigl(\{0\}\times[-1,1]\bigr).
  \]
  曲线 $G$ 的内部为空：任何开圆盘都含有不在这条曲线上的点。于是
  \[
    \overline C=\overline{\R^2\setminus G}
    =\R^2\setminus G^\circ=\R^2.
  \]
  同时
  \[
    C^\circ=\R^2\setminus\overline G
    =\R^2\setminus\left(G\cup(\{0\}\times[-1,1])\right).
  \]
  最后
  \[
    \partial C=\overline C\setminus C^\circ
    =G\cup\bigl(\{0\}\times[-1,1]\bigr).
  \]
\end{enumerate}

\problem{4}
\idea 对第一小题，考察指标 $m,n$ 是否趋于无穷；对第二小题，利用 $|\sin n|\leq1$。

\solution
\begin{enumerate}
  \item 令
  \[
    S=\left\{\frac1m+\frac1n:m,n\in\mathbb N\right\}.
  \]
  固定 $m$，令 $n\to\infty$，得到
  \[
    \frac1m+\frac1n\longrightarrow\frac1m.
  \]
  因而每个 $1/m$ 都是 $S$ 的极限点。再令 $m,n\to\infty$，可得 $0$ 也是极限点。

  下面说明没有其他极限点。设 $S$ 中一列互不相同的点
  \[
    s_k=\frac1{m_k}+\frac1{n_k}
  \]
  收敛。由于 $m,n$ 的地位对称，不妨交换每一项中的两个指标，使 $m_k\leq n_k$。若 $m_k$ 无界，可取子列使 $m_k\to\infty$；此时 $n_k\geq m_k$，故相应子列满足 $s_k\to0$。若 $m_k$ 有界，则可取无限子列使 $m_k$ 恒等于某个正整数 $m$。为了使这一子列中的 $s_k$ 互不相同，$n_k$ 必须取无穷多个值；再取子列可令 $n_k\to\infty$，于是 $s_k\to1/m$。原数列既然收敛，它的极限与任一收敛子列的极限相同，所以原极限只能是 $0$ 或某个 $1/m$。因此极限点集合恰为
  \[
    S'=\{0\}\cup\left\{\frac1m:m\in\mathbb N\right\}.
  \]

  \item 令
  \[
    T=\left\{\frac{\sin n}{n}:n\in\mathbb N\right\}.
  \]
  因为
  \[
    \left|\frac{\sin n}{n}\right|\leq\frac1n\longrightarrow0,
  \]
  所以这个数列趋于 $0$。这些项不可能从某一项起只取有限多个非零值，否则它们的绝对值会有一个正的下界，与趋于 $0$ 矛盾。因此存在互不相同的点趋于 $0$，故 $0$ 是极限点。

  对任何 $a\neq0$，取 $\varepsilon=|a|/3$。当 $n$ 充分大时，$|\sin n/n|<\varepsilon$，所以区间 $(a-\varepsilon,a+\varepsilon)$ 至多含有 $T$ 中有限多个点；再缩小邻域即可避开除 $a$ 外的这些有限个点。因此 $a$ 不是极限点。故
  \[
    T'=\{0\}.
  \]
\end{enumerate}

\problem{5}
\idea 要证 $x\in\overline{A\cap O}$，只需证明 $x$ 的每个开邻域都与 $A\cap O$ 相交。

\solution 任取 $x\in O$，并任取 $x$ 的开邻域 $U$。因为 $O$ 也是开集，$U\cap O$ 是 $x$ 的开邻域。又因为 $A$ 在 $X$ 中稠密，即 $\overline A=X$，每个非空开集都与 $A$ 相交，所以
\[
  (U\cap O)\cap A\neq\varnothing.
\]
等价地，$U\cap(A\cap O)\neq\varnothing$。这对 $x$ 的每个开邻域 $U$ 都成立，因此 $x\in\overline{A\cap O}$。由于 $x\in O$ 任意，得到
\[
  O\subseteq\overline{A\cap O}.
\]

\problem{6}
\idea 直接连续使用两次子空间拓扑的定义。

\solution 设 $W$ 是 $Z$ 中的任意开集。由于 $Z$ 是 $Y$ 的子空间，存在 $Y$ 中的开集 $V$，使得
\[
  W=Z\cap V.
\]
由于 $Y$ 是 $X$ 的子空间，又存在 $X$ 中的开集 $U$，使得
\[
  V=Y\cap U.
\]
注意 $Z\subseteq Y$，所以
\[
  W=Z\cap(Y\cap U)=(Z\cap Y)\cap U=Z\cap U.
\]
这说明 $Z$ 中每个开集都是某个 $X$ 中开集与 $Z$ 的交。反过来，任意 $X$ 中开集 $U$ 都满足
\[
  Z\cap U=Z\cap(Y\cap U),
\]
而 $Y\cap U$ 在 $Y$ 中开，所以 $Z\cap U$ 在 $Z$ 中开。因此 $Z$ 上原有的拓扑正是由 $X$ 诱导的子空间拓扑，即 $Z$ 是 $X$ 的子空间。

\problem{9}
\idea 分别处理“开”和“闭”两种情况，并使用子空间中开集、闭集的刻画。

\solution 若 $A$ 在 $Y$ 中开，则存在 $X$ 中的开集 $U$，使
\[
  A=Y\cap U.
\]
若 $Y$ 在 $X$ 中也开，那么 $A$ 是 $X$ 中两个开集的交，因此 $A$ 在 $X$ 中开。

若 $A$ 在 $Y$ 中闭，则存在 $X$ 中的闭集 $F$，使
\[
  A=Y\cap F.
\]
若 $Y$ 在 $X$ 中也闭，那么 $A$ 是 $X$ 中两个闭集的交，因此 $A$ 在 $X$ 中闭。这就分别证明了括号中的两个命题。

\problem{14}
\idea 证明同胚最方便的方法是找出显式逆函数，并验证原函数和逆函数都连续。

\solution 定义
\[
  h:\R\longrightarrow(0,1),\qquad h(x)=\frac{e^x}{1+e^x}.
\]
因为 $e^x>0$，所以 $0<h(x)<1$，函数确实取值于 $(0,1)$。若 $y=h(x)$，则
\[
  y=\frac{e^x}{1+e^x}
  \quad\Longleftrightarrow\quad
  y(1+e^x)=e^x
  \quad\Longleftrightarrow\quad
  e^x=\frac{y}{1-y}.
\]
当 $0<y<1$ 时，$y/(1-y)>0$，因而可以取对数，得到唯一的
\[
  x=\log\frac{y}{1-y}.
\]
因此 $h$ 是双射，其逆函数为
\[
  h^{-1}(y)=\log\frac{y}{1-y},\qquad 0<y<1.
\]
$h$ 由指数函数及四则运算复合而成，且分母始终不为零，所以在 $\R$ 上连续。$h^{-1}$ 由连续函数 $y$、$1-y$、除法和对数复合而成，在 $(0,1)$ 上也连续。因此 $h$ 是从 $\R$ 到 $(0,1)$ 的同胚。

\problem{16}
\idea 用任意子集的示性函数，把“每个实值函数都连续”转化为“每个子集都开”。

\solution 任取 $A\subseteq X$，定义示性函数（indicator function）
\[
  \chi_A:X\longrightarrow\R,
  \qquad
  \chi_A(x)=
  \begin{cases}
    1,&x\in A,\\
    0,&x\notin A.
  \end{cases}
\]
由题设，$\chi_A$ 连续。实数直线中的区间 $(1/2,3/2)$ 是开集，而
\[
  \chi_A^{-1}\bigl((1/2,3/2)\bigr)=A.
\]
连续映射把开集的原像变成开集，所以 $A$ 在 $X$ 中开。由于 $A$ 是任意子集，$X$ 的每个子集都是开集。因此 $X$ 必须赋予离散拓扑（discrete topology）
\[
  \mathcal T=\mathcal P(X).
\]
反过来，在离散空间上任意映射都是连续的，所以这个条件也确实充分。

\problem{17}
\idea 连续性检查开集的原像；不是同胚则找一个通常开集，说明其像在余有限拓扑中不开。

\solution 设定义域 $\R$ 取通常拓扑，值域 $X$ 的开集为 $\varnothing$ 以及补集为有限集的子集。映射
\[
  f:\R\longrightarrow X,
  \qquad f(x)=x
\]
在集合意义下是恒等映射，因而是双射。

任取 $X$ 中的开集 $U$。若 $U=\varnothing$，则 $f^{-1}(U)=\varnothing$，在 $\R$ 中开。若 $U\neq\varnothing$，则存在有限集 $F\subseteq\R$，使
\[
  U=\R\setminus F.
\]
有限集 $F$ 在通常实数直线中闭，所以 $\R\setminus F$ 在通常拓扑中开。又因为 $f$ 是恒等映射，
\[
  f^{-1}(U)=U=\R\setminus F.
\]
因此所有 $X$ 中开集的原像都在通常 $\R$ 中开，故 $f$ 连续。

但是 $f$ 不是同胚。例如 $(0,1)$ 在通常 $\R$ 中开，而
\[
  f((0,1))=(0,1)
\]
在 $X$ 中不开，因为它的补集 $(-\infty,0]\cup[1,+\infty)$ 不是有限集。因此 $f$ 不是开映射；等价地，逆映射 $f^{-1}:X\to\R$ 不连续。故 $f$ 虽是连续双射，却不是同胚。

\vfill
\noindent{\small\textbf{题目来源：}M. A. Armstrong, \textit{Basic Topology}，第 2 章，第 1、2、3、4、5、6、9、14、16、17 题。}

\end{document}
